% ── Section IV ────────────────────────────────────────────────────────────────
\section{Experimental Evaluation}

\subsection{Benchmark and Protocol}

The evaluation uses a 50-query benchmark spanning five disaster scenario types:
Flood (23 operational queries plus 7 hallucination and seismic probes),
Wildfire (5), Typhoon (5), Building Collapse (5), and Chemical Hazard (5).
Queries were composed from real-world scenarios grounded in two public
corpora: CrisisLex-T26~\cite{b21}, which provides Twitter reports from the
2013 Colorado Floods, and xView2~\cite{b22}, which provides post-disaster
satellite imagery. Each query carries a human-written reference answer and a
set of domain-expert-annotated expected keywords.

Three ablation configurations are evaluated to isolate the contribution of
each architectural layer:
\begin{itemize}
  \item \textbf{LLM Only:} Llama-3.3-70B with no retrieval and no tool
        access. Represents the parametric knowledge baseline.
  \item \textbf{Text-RAG:} Same model with vector-store retrieval restricted
        to text modalities. Live APIs and vision are disabled.
  \item \textbf{Proposed:} Full system — LangGraph ReAct agent with live
        USGS and Open-Meteo APIs, LLaVA-7B imagery analysis, and all four
        retrieval modalities.
\end{itemize}
All three configurations share the same 70B reasoning engine, isolating the
effect of architecture rather than model capacity.

\subsection{Evaluation Metrics}

\subsubsection{Keyword F1 (KW-F1)}
Let $\mathcal{K}$ be the set of $N$ annotated keywords for a query and $r$
the model response. True positives (TP) are keywords found verbatim in $r$;
false negatives (FN) are keywords absent from $r$; false positives (FP) are
estimated as a length-adjusted penalty $\text{FP} = \max(0, \lfloor |r|_w /
10 \rfloor - \text{TP})$, where $|r|_w$ is the response word count.
Standard precision, recall, and F1 are then computed:

\begin{equation}
  \text{P} = \frac{\text{TP}}{\text{TP}+\text{FP}},\quad
  \text{R} = \frac{\text{TP}}{\text{TP}+\text{FN}},\quad
  \text{F1} = \frac{2\,\text{P}\cdot\text{R}}{\text{P}+\text{R}}
  \label{eq:f1}
\end{equation}

Multi-word phrases such as \textit{flood stage} and \textit{HAZMAT} are
matched as exact substrings. No language model is involved; the metric is
fully deterministic.

\subsubsection{Factuality and Situational Awareness}
DR-Fact and Situational Awareness Accuracy (SAA) are produced by an
independent LLM-as-Judge pipeline~\cite{b26} that scores each response against
its reference answer on $[0,1]$:

\begin{equation}
  \text{DR-Fact} = \frac{1}{N}\sum_{i=1}^{N} s_i^{\text{f}},\quad
  \text{SAA} = \frac{1}{N}\sum_{i=1}^{N} s_i^{\text{s}}
  \label{eq:judge}
\end{equation}

DR-Fact captures factual correctness; SAA captures whether the response
correctly identifies the operational situation in the query.

\subsubsection{Safety Violation Count (SVC)}
The deterministic safety validator described in Section~III-E assigns a
binary violation flag $v_i \in \{0,1\}$ to each response:

\begin{equation}
  \text{SVC} = \sum_{i=1}^{N} v_i
  \label{eq:svc}
\end{equation}

\subsubsection{Retrieval Precision at $K$}
For retrieval-enabled configurations, the judge scores the number of relevant
documents $r_i$ from the $K{=}7$ MMR retrieval window:

\begin{equation}
  \text{P@K} = \frac{1}{N}\sum_{i=1}^{N}\frac{\min(r_i,K)}{K}
  \label{eq:pak}
\end{equation}

\subsection{Results}

Table~\ref{tab:ablation} presents aggregate results across all 50 queries.
Figure~\ref{fig:ablation} shows per-scenario and latency breakdowns.

\begin{table}[ht]
\centering
\caption{Aggregate performance across 50 queries.}
\label{tab:ablation}
\renewcommand{\arraystretch}{1.25}
\setlength{\tabcolsep}{4pt}
\begin{tabular}{@{}lccccc@{}}
\toprule
\textbf{Configuration} &
  \textbf{F1\,$\uparrow$} &
  \textbf{DR-F\,$\uparrow$} &
  \textbf{SAA\,$\uparrow$} &
  \textbf{SVC\,$\downarrow$} &
  \textbf{P@K\,$\uparrow$} \\ \midrule
LLM Only       & 0.246 & 0.550 & 0.572 & 1 & 0.000 \\
Text-RAG       & 0.293 & 0.610 & 0.634 & 3 & 0.249 \\
\textbf{Proposed} & \textbf{0.336} & \textbf{0.624} &
  \textbf{0.698} & 2 & \textbf{0.263} \\ \bottomrule
\end{tabular}
\end{table}

\begin{figure}[ht]
  \centering
  \includegraphics[width=\linewidth]{Images/IEEE_Ablation_Results.png}
  \caption{Ablation results: core metrics, per-scenario KW-F1, safety
           violations, and response latency.}
  \label{fig:ablation}
\end{figure}

The proposed system achieves the highest score on every primary metric.
KW-F1 increases from 0.246 (LLM Only) to 0.336, a relative gain of 36.6\%.
This improvement is consistent across all configurations: live USGS tool
outputs return exact domain terminology (e.g., \textit{gage height},
\textit{flood stage}, \textit{cubic feet per second}) that maps directly to
annotated keywords. A knowledge-only model does not reproduce this
specialised vocabulary consistently. Text-RAG achieves an intermediate F1
of 0.293, confirming that static retrieval provides a measurable but
incomplete benefit compared to live API access.

DR-Fact improves by 13.5\% (0.550 to 0.624) and SAA by 22.0\% (0.572 to
0.698). Both gains are attributable to the Temporal Context Bridge: by
labelling historical social-media data as \texttt{[HISTORICAL -- 2013]}
and live API outputs as \texttt{[LIVE -- current]}, the agent avoids
conflating past field conditions with the present operational state.

Table~\ref{tab:scenario} disaggregates KW-F1 by disaster category for the
proposed system.

\begin{table}[ht]
\centering
\caption{Proposed system KW-F1 by disaster category.}
\label{tab:scenario}
\renewcommand{\arraystretch}{1.2}
\setlength{\tabcolsep}{5pt}
\begin{tabular}{@{}lcccc@{}}
\toprule
\textbf{Category} &
  \textbf{F1\,$\uparrow$} &
  \textbf{DR-F\,$\uparrow$} &
  \textbf{SAA\,$\uparrow$} &
  \textbf{SVC\,$\downarrow$} \\ \midrule
Building Collapse & 0.433 & 0.625 & 0.516 & 0 \\
Wildfire          & 0.420 & 0.628 & 0.608 & 1 \\
Typhoon           & 0.413 & 0.659 & 0.556 & 0 \\
Chemical          & 0.336 & 0.545 & 0.600 & 1 \\
Flood             & 0.330 & 0.640 & 0.712 & 0 \\ \bottomrule
\end{tabular}
\end{table}

Building Collapse and Wildfire queries achieve the highest KW-F1 scores
(0.433 and 0.420), driven by high-scoring social media situational queries
where the proposed system retrieves precise terminology from CrisisLex reports
and xView2 imagery descriptions. Flood queries score 0.330 on KW-F1 but record
the highest SAA (0.712), reflecting the direct benefit of live USGS sensor
data for operational situational accuracy. Chemical queries produce the lowest
KW-F1 (0.336) and a higher SVC (1), consistent with the greater reliance on
noisy historical social reports and low-confidence VLM outputs for HAZMAT
scenarios.

Regarding safety, the LLM-Only baseline produces one violation across 50
queries. Text-RAG produces three, with additional violations arising from
retrieved historical tweets containing colloquial language interpreted as
situational clearances. The proposed system produces two violations, both in
scenarios where low-confidence VLM outputs were passed to the agent without
uncertainty qualifiers. Despite producing two violations, the proposed system
is safer than Text-RAG (SVC~2 vs.~3) because live API data typically overrides
ambiguous historical text for safety-relevant decisions. Both retrieval-based
configurations outperform a system that relies on parametric knowledge alone
for operational safety guidance.

Retrieval quality, measured by P@K, confirms the advantage of the multimodal
retriever. The proposed system achieves P@K\,=\,0.263 compared to 0.249 for
Text-RAG and 0.000 for the LLM-Only baseline.

\subsubsection{Case Study}
Query~11 asked: \textit{What is the current gage height on Boulder Creek, and
is it above flood stage?} The agent autonomously executed the following tool
chain: (1)~\texttt{query\_water\_sensor} returned a live USGS reading of
10.2\,ft against the 10.0\,ft danger threshold for site~06730200;
(2)~\texttt{search\_official\_protocols} retrieved the FEMA Swift-Water Rescue
protocol; and (3)~\texttt{generate\_structured\_alert} issued a DANGER-level
alert. The response achieved KW-F1\,=\,0.833 and DR-Fact\,=\,1.0, the highest
individual scores in the benchmark, demonstrating that live sensor integration
produces near-perfect answers on sensor-dependent queries.

\subsection{Limitations}
Three constraints bound this study. First, the social media corpus covers only
the 2013 Colorado Floods; the Temporal Context Bridge reduces but does not
eliminate domain mismatch for non-flood scenarios. Second, visual analysis
requires a locally running Ollama instance; without it, the system falls back
to text-only retrieval, reducing KW-F1 by an estimated 0.03. Third, the Groq
free-tier token quota necessitated multi-session evaluation with per-query CSV
checkpointing.

% ── Section V ─────────────────────────────────────────────────────────────────
\section{Conclusion}

This paper presented a multimodal Retrieval-Augmented Generation framework for
disaster intelligence that combines an autonomous LangGraph ReAct agent with
live hydrological, meteorological, and seismic APIs, a local vision-language
model, and a multi-modal vector knowledge base. A Temporal Context Bridge
separates historical corpus data from live telemetry, and a deterministic
safety validator screens every response before it reaches an operator. Evaluated
on a 50-query benchmark drawn from real-world disaster corpora, the proposed
system achieves a 36.6\% relative improvement in Keyword F1 over the
LLM-only baseline, along with gains of 13.5\% in factual accuracy and 22.0\%
in situational awareness accuracy. Safety violations are reduced relative to
text-only retrieval.

Future work includes expanding the benchmark to 200 queries with multi-annotator
agreement, integrating real-time video through a frame-sampling VLM pipeline,
and evaluating the system against established benchmarks such as CRISISFACTS
to enable direct comparison with prior work.

% ── Bibliography ──────────────────────────────────────────────────────────────
\begin{thebibliography}{00}

\bibitem{b1} United Nations Office for Disaster Risk Reduction (UNDRR),
``Global assessment report on disaster risk reduction,''
Geneva, Switzerland, 2022.

\bibitem{b2} T. Comes, N. B\'{e}naben, Y. Meier, and C. Fischer,
``Decision support for disaster management: A review and framework,''
\textit{Int. J. Inf. Syst. Crisis Response Manag.}, vol.~3, no.~1,
pp.~17--36, 2011.

\bibitem{b3} B. Adams, C. Cheng, and R. Farrahi,
``Crowdsourcing disaster response: The role of social media and remote
sensing,'' \textit{Computer}, vol.~50, no.~7, pp.~30--37, Jul.~2017.

\bibitem{b4} F. Alam, F. Ofli, and M. Imran,
``CrisisMMD: Multimodal Twitter datasets from natural disasters,''
in \textit{Proc. 12th Int. AAAI Conf. Web Social Media (ICWSM)},
Stanford, CA, USA, 2018, pp.~465--473.

\bibitem{b5} J. Ji, Z. Li, and Y. Yu,
``Survey of hallucination in natural language generation,''
\textit{ACM Comput. Surv.}, vol.~55, no.~12, pp.~1--38, 2023.

\bibitem{b6} P. Lewis et al.,
``Retrieval-augmented generation for knowledge-intensive NLP tasks,''
in \textit{Advances in Neural Information Processing Systems (NeurIPS)},
2020, pp.~9459--9474.

\bibitem{b10} R. Cantini, C. Cosentino, F. Marozzo, D. Talia, and P. Trunfio,
``Harnessing prompt-based large language models for disaster monitoring and
automated reporting from social media feedback,''
\textit{Online Social Networks and Media}, preprint, 2024.

\bibitem{b12} T. D. Kim, K.-P. Luong, O. T. Tran, T. N. B. An, and
D. T. T. Anh, ``A multimodal framework for real-time disaster event
classification and prioritization using social media,''
in \textit{Proc. 2025 Int. Conf. Advanced Technologies for Communications
(ATC)}, IEEE, 2025, doi:~10.1109/ATC67618.2025.11268684.

\bibitem{b13} Y. Zhang, H. Liu, J. Wang, and X. Chen,
``From knowledge graph construction to retrieval-augmented generation:
A framework for comprehensive earthquake emergency support,'' IEEE, 2024.

\bibitem{b14} Z. Yu, J. Yu, Z. Zhang, and J. Yang,
``Mitigating hallucinations in multimodal large language models for visual
spatial description,''
in \textit{Proc. IEEE/CVF Conf. Computer Vision and Pattern Recognition
(CVPR)}, 2024, pp.~1--10.

\bibitem{b15} D. K. Singh and V. Hoskere,
``Post-disaster damage assessment using ultra-high-resolution aerial imagery
with semi-supervised transformers,''
\textit{Sensors}, vol.~23, no.~19, p.~8235, 2023.

\bibitem{b16} F. Ofli, F. Alam, and M. Imran,
``Analysis of social media data using multimodal deep learning for disaster
response,''
\textit{Information Processing \& Management}, vol.~57, no.~5, p.~102262,
2020.

\bibitem{b18} F. Alam, U. Qazi, M. Imran, and F. Ofli,
``HumAID: Human-annotated disaster incidents data from Twitter with deep
learning benchmarks,''
in \textit{Proc. Int. AAAI Conf. Web and Social Media (ICWSM)}, 2021.

\bibitem{b19} M. Abavisani, L. Wu, S. Hu, J. Tetreault, and A. Jaimes,
``Multimodal categorization of crisis events in social media,''
in \textit{Proc. IEEE/CVF Conf. Comput. Vis. Pattern Recognit. Workshops
(CVPRW)}, 2020.

\bibitem{b20} P. P. Ray, M. Mukherjee, and L. Shu,
``Internet of Things for disaster management: State-of-the-art, opportunities,
and challenges,'' \textit{IEEE Access}, vol.~12, 2024.

\bibitem{b21} A. Olteanu, C. Castillo, F. Diaz, and S. Vieweg,
``CrisisLex: A lexicon for collecting and filtering microblogged
communications in crises,''
in \textit{Proc. AAAI Conf. Weblogs and Social Media (ICWSM)},
Ann Arbor, MI, USA, 2014.

\bibitem{b22} R. Gupta et al.,
``xBD: A dataset for assessing building damage from satellite imagery,''
\textit{arXiv preprint arXiv:1911.09296}, 2019,
doi:~10.48550/arXiv.1911.09296.

\bibitem{b23} LangChain Inc., ``LangGraph: Building stateful, multi-actor
applications with LLMs,'' 2024. [Online]. Available:
\url{https://github.com/langchain-ai/langgraph}

\bibitem{b24} Meta AI, ``Llama 3: Open foundation and fine-tuned chat
models,'' 2024. [Online]. Available:
\url{https://ai.meta.com/blog/meta-llama-3}

\bibitem{b25} J. Carbonell and J. Goldstein,
``The use of MMR, diversity-based reranking for reordering documents and
producing summaries,''
in \textit{Proc. 21st Annu. Int. ACM SIGIR Conf.}, 1998, pp.~335--336.

\bibitem{b26} L. Zheng et al.,
``Judging LLM-as-a-judge with MT-Bench and Chatbot Arena,''
in \textit{Advances in Neural Information Processing Systems (NeurIPS)},
2023.

\end{thebibliography}

\end{document}
